The Iron Throne had never once, in all the years she had studied it from every angle a woman was permitted to study anything from, looked comfortable to sit upon, and Cersei had come to understand, this past year, that this had always been rather the point of it.

She did not sit it often now. The babe had made the climb up its black iron steps a labor she preferred to spare herself when there was no audience worth impressing, and there was, this morning, no audience at all — only Qyburn, waiting at the base of the dais with the particular stillness of a man who had learned long ago that a queen’s silences were rarely empty ones and were best not interrupted before they’d finished with themselves.

“Tell me the numbers again,” she said, when she judged the silence had done its work.

“The Golden Company sails with the tide favoring them, by Lord Greyjoy’s last word — two days out from the Blackwater, perhaps three, if the winter seas trouble them.” Qyburn’s voice remained the same mild, papery register it always did, whether he was reporting a fleet’s position or describing, as he sometimes still did unbidden, the finer details of whatever unfortunate creature his experiments had most recently required. “Ten thousand sellswords, elephants included, though I am given to understand the beasts have not traveled the crossing well.”

“And the city garrison?”

“Sufficient to hold the walls against any force presently marching against them, Your Grace. Insufficient, I would judge, against certain other possibilities we have discussed, should they prove more than a maester’s fireside tale.”

She did not ask him to elaborate on which possibilities he meant. They had discussed the dead enough times now that naming them again felt less like prudence and more like inviting something to overhear its own name spoken and take an interest. She had never been a superstitious woman — had made something of a private point, most of her life, of despising the superstitions that ruled weaker minds around her, Robert’s septons and Ned Stark’s gods and her own father’s blind faith in the family name outlasting any single member of it — and yet she found, this past year, an old prophecy returning to her at hours she did not choose, unbidden as a draft finding its way beneath a door someone had sworn was sealed.

\textit{Gold shall be their crowns, and gold their shrouds.} Maggy the Frog had promised her that, in a tent that smelled of blood and goat, when Cersei had been a girl of ten who believed the whole exercise a game. \textit{And when your tears have drowned you, the valonqar shall wrap his hands about your pale white throat and choke the life from you.}

She had spent the better part of thirty years deciding the little brother in question was Tyrion, because it was easier, infinitely easier, to fear a hatred she had already earned than to examine the alternative — that \textit{valonqar} meant nothing more specific than \textit{little brother}, in the Old Tongue, and that she had two of them, and had always, if she was entirely honest with the parts of herself she rarely let speak, suspected which one the prophecy actually meant and simply refused the thought each time it surfaced, the way a tongue refuses to press against a rotting tooth.

She did not let herself finish that thought now either. She had a child coming — his child, hers, theirs, the only true thing left standing between her and a world that had taken every other true thing she’d ever had reason to trust — and prophecy, she reminded herself, as she had reminded herself a thousand times before, was a weapon wielded far more often by the fear it created in the hearer than by any power in the words themselves. Maggy the Frog was three decades dead. Cersei Lannister was not.

“There is one more matter,” Qyburn said, when she had been quiet long enough that he judged it safe to raise it. “Ser Jaime.”

Her hand went, without her quite willing it, to the small swell beneath her gown, still modest enough to hide beneath good tailoring, though not for very much longer. “What of him.”

“No further word since he crossed the Neck, Your Grace. My little birds lose him past there more often than not — the North keeps its secrets better than the South does, and keeps them from paid ears especially. I expect he has reached Winterfell by now, or very nearly.”

“Do you.” She said it flatly, and did not ask the question underneath it, the one she had not asked anyone since the morning she’d told Jaime plainly that she would not beg him to stay and had meant it, right up until the door closed behind him and she had discovered, alone in a room too large for one woman and an unborn child, that meaning a thing and surviving the meaning of it were two entirely different exercises.

She had told him Winterfell’s war was no war of theirs. She had believed it, mostly, the way she believed most things that were more convenient than examined. It was easier to disbelieve an army of the dead than to admit that her own throne — this one, black and ugly and never once comfortable, precisely as it had never been designed to be — meant less to her, in the smallest hours of the night, than the question of whether the father of her last living child would come home to see it born.

She did not say any of that to Qyburn, who would have recorded it somewhere useful and reported it somewhere more useful still, and who she trusted for his competence rather than his discretion, a distinction she had learned to keep very carefully separated in her own mind.

“Robert Strong,” she said instead, turning the conversation the way she turned most conversations that threatened to go somewhere she had not chosen. “Is he fit to stand at the Red Keep’s gates, should the Golden Company arrive before I’ve had proper word of them?”

“He does not tire, Your Grace. He does not fear. He does not, so far as I have been able to determine, do very much of anything at all beyond precisely what he is told.” Qyburn’s small smile was the closest thing to pride she had ever seen him permit himself. “In that respect he is, I would humbly submit, the most reliable servant currently in your employ.”

“Small comfort, that the dead make better subjects than the living.”

“I have found it so myself, Your Grace, on more occasions than I would care to enumerate.”

“There is a third matter,” she said, before he could take his leave, “and I want it settled before the Golden Company’s sails are close enough to count. The wildfire stores beneath the city — what remains of them, and what more can be made, if more is needed.”

“Considerable stores remain, Your Grace, from what the Mad King’s alchemists laid down and what your own reign has added since. Whether more can be \textit{made}, in the time we may have—“ Qyburn hesitated, an unusual thing for him. “That knowledge does not live in a ledger. It lives, so far as I am aware, in one very old man’s skull, and I am given to understand the skull in question is failing rather faster than the war is arriving.”

“Wisdom Hallyne.”

“You remember him.”

“I remember every man in this city who has ever been useful to me, Lord Qyburn. It is a shorter list than you might think, which is precisely why I remember it so well.” She was already rising, the babe’s weight making the motion slower than it once would have been, slower than she liked anyone to see, though Qyburn had long since earned the privilege of seeing her at less than her best. “Take me to him.”

\sceneBreak

The Guildhall of the Alchemists smelled, even from the street, of something both sweetly rotten and faintly metallic, a smell Cersei remembered from girlhood without ever having troubled to learn its cause, and Hallyne’s sickroom, three flights up a stair no manservant had thought to make easier for a dying man to climb, smelled of it worst of all — as though the substance the old pyromancer had spent his whole life coaxing into being had, in the end, decided to take up residence inside him rather than release him.

He was smaller than she remembered, folded into a bed too large for what remained of him, and his eyes, when he opened them to her, had gone the same sickly green-white glow she recalled from the vaults beneath the Great Sept all those years ago — not a trick of firelight, she understood now, watching them catch what little sun came through the shutters, but something the wildfire itself had left behind in him, the way a lifetime spent near a poison eventually becomes a kind of poisoning.

“Your Grace.” His voice was a dry rattle, delighted despite itself. “I did not think to be honored twice in one lifetime.”

“I need the substance, Wisdom Hallyne. In greater quantity than this city has ever held at once, and I need it before winter finishes what your own guild started centuries ago.” She did not soften the request, because she had learned, watching her father conduct a thousand such transactions, that dying men respected plain need far more than they respected false comfort, the same lesson the North’s lords had apparently taught Sansa Stark, though Cersei would sooner choke than credit a Stark with anything worth learning from. “Tell me what remains to be told, and tell it to Lord Qyburn plainly enough that the knowledge survives you.”

Hallyne’s ruined eyes drifted, for a moment, somewhere past her shoulder, to some vault of memory she had no wish to follow him into. “It does not like to be rushed, Your Grace. The substance. It has its own patience, its own... appetite, one might say, and a man who forgets that in his haste loses considerably more than time for the forgetting.” A wet, rattling breath. “But there are ways to hasten what cannot be hurried. Older vats. Deeper cellars. My apprentice Garth knows the proportions near as well as I do myself — better, perhaps, in years to come, once I am past knowing anything at all.”

“Then Garth will answer to Lord Qyburn from this day forward,” Cersei said, “and you, Wisdom Hallyne, will spend whatever days remain to you making certain he has forgotten nothing worth remembering.”

The old man’s laugh dissolved into coughing, and when it passed he looked at her with something that might, in a less ruined face, have been respect. “They say wildfire wants to be free, Your Grace. That is the old jape among my order, the first thing a novice is taught to fear. I have spent sixty years keeping it from getting what it wants.” His green eyes found hers, unblinking. “I wonder, sometimes, whether that was ever truly a kindness to this city.”

Cersei did not answer that. She had long since stopped believing questions of kindness had any bearing on what a city needed to survive, and she left him to Qyburn’s attending and the Guildhall’s sour, metallic air, already turning her thoughts toward vats and cellars and a substance patient enough to have waited centuries for its freedom, and willing, this once, to let it wait a little longer still, so long as it waited on her terms and no one else’s.

She rode back to the Red Keep in a litter she detested and would not admit to needing, and sat a while longer alone beneath the Iron Throne rather than upon it, one hand resting where the babe turned inside her — his child, Jaime’s, the last true Lannister gold left to matter once every other claim to that name had burned or bent the knee or ridden north to die for people who had never once, in three hundred years, thanked a Lannister for anything — and thought, not for the first time, that power had always been a simpler thing to hold than to keep, and that she had spent her whole life mistaking the first for the second, and that she did not, any longer, entirely believe there was time left to learn the difference before someone came to collect what she owed for the mistake.

Outside, beyond the Red Keep’s walls, the city went about its business in the thin grey light of a winter that had lasted longer than any within living memory, unaware, mostly, of dragons or dead men or debts owed to old prophecies in old tents that smelled of blood and goat — and Cersei Lannister sat beneath a throne built from the melted swords of a war fought before her grandfather’s grandfather had been born, and waited, as she had always waited, to see which of the many things she feared arrived first to collect.